% !TEX root = ../main.tex
\chapter{LaTeX examples}
\label{app:examples}
This appendix demonstrates the scaffold's conventions. It is illustrative
material, not a claim about the thesis research. Remove it when writing.

\section{Equations and definitions}
Let \(x\in\R^n\) be an unknown vector and
\(p\in\mathcal P\subseteq\R^m\) a parameter vector. Consider
\begin{equation}
  F(x,p)=0,\qquad F:\R^n\times\mathcal P\longrightarrow\R^n.
  \label{eq:residual}
\end{equation}
The \gls{residual} \(F\) measures the mismatch in the equations.
Its \gls{jacobian} with respect to \(x\) is written \(F_x\).
\Cref{eq:residual} defines an implicit problem.

\begin{definition}[Euclidean norm]
  \label{def:norm}
  For \(x\in\R^n\), the Euclidean norm is
  \[
    \norm{x}=\sqrt{\inner{x}{x}}=\left(\sum_{i=1}^{n}x_i^2\right)^{1/2}.
  \]
\end{definition}

\begin{proposition}[Zero norm]
  \label{prop:zero-norm}
  For \(x\in\R^n\), \(\norm{x}=0\) if and only if \(x=0\).
\end{proposition}
\begin{proof}
  If \(\norm{x}=0\), then \(\sum_i x_i^2=0\).
  Every summand is nonnegative, so each \(x_i^2=0\) and hence \(x=0\).
  Conversely, substitution of \(x=0\) in \cref{def:norm} gives \(\norm{x}=0\).
\end{proof}

\begin{remark}
  \label{rem:scaling}
  A small residual alone does not guarantee a small error in the unknown.
  The relationship depends on conditioning and the scales of the equations.
\end{remark}
The environments share a chapter counter; \cref{prop:zero-norm,rem:scaling}
illustrate references to different statement types.

\section{Acronyms and citations}
A time-dependent model may take the form of an \gls{ode} or a \gls{dae}.
Derivatives of computational expressions can be evaluated using \gls{ad}.
Later uses remain abbreviated: \gls{ode}, \gls{dae}, and \gls{ad}.
Definitions still belong in the prose, with the lists serving as reference aids.

For a citation example, \textcite{latex-project} describes the LaTeX project.
Bibliography formatting is provided by \texttt{biblatex}\autocite{biblatex}.
These documentation references are sample entries to replace with your literature.

\section{An algorithm}
For fixed \(p\), \cref{alg:newton} illustrates a Newton iteration.
It assumes each required Jacobian is nonsingular. The returned Boolean records
whether the requested residual tolerance was achieved; it is not an error bound.

\begin{algorithm}[htbp]
  \caption{Newton iteration for a fixed parameter vector}
  \label{alg:newton}
  \begin{algorithmic}[1]
    \Require Initial guess \(x_0\), parameter \(p\), tolerance \(\varepsilon>0\), iteration limit \(K\geq1\)
    \State \(x\gets x_0\)
    \For{\(k=0,\ldots,K-1\)}
    \If{\(\norm{F(x,p)}\leq\varepsilon\)}
    \State \Return \((x,\mathrm{true})\)
    \EndIf
    \State Solve \(F_x(x,p)s=-F(x,p)\) for \(s\)
    \State \(x\gets x+s\)
    \EndFor
    \State \Return \((x,\norm{F(x,p)}\leq\varepsilon)\)
  \end{algorithmic}
\end{algorithm}

\section{Tables, figures and code}
The numbers in \cref{tab:example} are illustrative formatting data.
The table uses horizontal rules and decimal alignment.
\begin{table}[htbp]
  \centering
  \caption{Illustrative numerical values, not experimental results.}
  \label{tab:example}
  \begin{tabular}{@{}l S[table-format=1.3] S[table-format=2.1]@{}}
    \toprule
    Case & {Residual norm} & {Time (\unit{\milli\second})} \\
    \midrule
    A    & 0.125           & 12.4                          \\
    B    & 0.031           & 18.7                          \\
    \bottomrule
  \end{tabular}
\end{table}

For a physical quantity, write \(t=\qty{2.5}{\milli\second}\).
Figure files belong in \texttt{figures/}; the source below is a pattern to use
once a real figure exists.
\begin{lstlisting}[language=TeX]
\begin{figure}[htbp]
  \centering
  \includegraphics[width=0.8\linewidth]{convergence.pdf}
  \caption{Explain what is plotted and why it matters.}
  \label{fig:convergence}
\end{figure}
\end{lstlisting}
A short source excerpt can be included directly:
\begin{lstlisting}[language=Python,caption={A scalar residual used only as a syntax example.},label={lst:residual}]
def residual(x, p):
    return x*x - p
\end{lstlisting}
\Cref{lst:residual} represents \(F(x,p)=x^2-p\).
